\section{Discussion}

In this work, we have introduced \textbf{Active-Bidder}, the first framework to reformulate real-time bidding (RTB) through the lens of Active Inference and the Free Energy Principle (FEP). By shifting the objective from myopic reward maximization (e.g., CTR or GMV) to the minimization of Variational Free Energy (VFE), we have demonstrated that advertising agents can maintain ``user homeostasis''---balancing exploratory information gain with exploitative preference fulfillment. This section contextualizes our findings within the broader landscape of AI safety and large-scale recommender systems.

\subsection{Theoretical Significance: Beyond Myopic Exploitation}
Traditional reinforcement learning (RL) approaches in computational advertising typically treat the user as a passive stochastic environment that emits rewards \cite{sutton2018reinforcement}. This perspective often leads to the ``filter bubble'' phenomenon and the optimization of addictive feedback loops, as agents exploit high-variance behaviors to maximize short-term clicks \cite{milano2020recommender}. In contrast, Active-Bidder treats the user as a \textit{self-evidencing agent} \cite{friston2010action}. By minimizing VFE, our agent inherently accounts for the epistemic value of an ad---asking not just ``will they click?'' but ``will this ad help me better understand the user's latent state?'' This aligns with recent shifts toward ``Human-Centric AI'' that prioritize long-term user agency over short-term conversion \cite{shneiderman2022human, zhang2024trustworthy}.

\subsection{Limitations: The Computational Cost of Inference}
Despite its theoretical elegance, the primary bottleneck for the deployment of Active-Bidder is the computational complexity associated with calculating the Variational Free Energy in real-time. RTB environments typically demand latencies below 50ms. The current implementation of Active-Bidder relies on amortized variational inference \cite{kingma2013auto}, which offloads much of the computational burden to an offline training phase. However, the E-step required to update the agent's posterior beliefs about the user's state ($\mathbf{s}_t$) still necessitates a forward pass through a deep generative model.

Furthermore, calculating the \textit{Expected Free Energy} (EFE) for future policies involves a summation over all possible future observations ($\mathbf{o}_{\tau}$), which scales exponentially with the planning horizon. While we utilized a mean-field approximation to simplify these dependencies, this can lead to an underestimation of the ``epistemic surprise'' in complex, multi-modal user distributions \cite{buckley2017free, smith2022active}. Future implementations may require hardware-accelerated predictive coding architectures or specialized neuromorphic chips to achieve the necessary throughput for global-scale ad exchanges \cite{scellier2017equilibrium, davies2024neuromorphic}.

\subsection{Societal Impact: Combating Clickbait and Manipulation}
The ethical implications of Active-Bidder are significant. Most current advertising algorithms are optimized for ``engagement,'' a metric that is easily ``hacked'' by clickbait or emotionally manipulative content that triggers dopamine-driven responses \cite{seaver2019algorithms}. Because Active-Bidder seeks to minimize the discrepancy between the user's internal generative model and the external stimuli (ads), it naturally penalizes ``surprising'' but ``uninformative'' content. 

An ad that promises one thing (high click-through probability) but delivers something irrelevant to the user’s long-term preferences (high VFE) will be discarded by the agent as it attempts to maintain the user's state of ``low surprise'' or homeostasis \cite{constant2022biosemiotics}. Consequently, Active-Bidder provides a mathematical foundation for \textit{alignment-by-design} in recommender systems. By reducing the reliance on high-entropy, manipulative stimuli, we can foster a digital ecosystem that respects the cognitive boundaries of the user \cite{pariser2011filter, hancox2024digital}.

\subsection{Future Work: Multi-Agent Active Inference}
The current iteration of Active-Bidder assumes a stationary or slowly evolving user model. A more realistic, albeit more complex, formulation involves a \textit{Multi-Agent Active Inference} (MA-AI) framework \cite{albarracin2023epistemic}. In this scenario, both the Advertiser and the User are modeled as Active Inference agents engaged in a ``communication game.'' The advertiser minimizes VFE by providing relevant information, while the user minimizes VFE by selecting which ads to attend to.

This co-evolutionary perspective opens several avenues for research:
\begin{enumerate}
    \item \textbf{Shared Generative Models}: Investigating how a ``shared world model'' emerges between a platform and its users through repeated interaction, potentially leading to more efficient ``zero-shot'' ad delivery \cite{friston2024morphogenesis}.
    \item \textbf{Privacy-Preserving Inference}: Using federated Active Inference to update user models locally on-device, ensuring that the latent states ($\mathbf{s}_t$) representing sensitive user preferences never leave the local environment, thus adhering to strict GDPR and CCPA standards \cite{tan2022federated, li2025privacy}.
    \item \textbf{Dynamic Auctions}: Extending the VFE objective to the auction mechanism itself, where the exchange seeks to minimize the ``Total Free Energy'' of the ecosystem rather than just maximizing revenue.
\end{enumerate}

\subsection{Conclusion}
Active-Bidder represents a departure from the ``click-maximization'' paradigm that has dominated the last decade of ad-tech. By grounding bidding strategies in the first principles of biological self-organization, we have shown that it is possible to create advertising agents that are both effective and ethically aligned. As we move toward an era of increasingly autonomous AI agents, frameworks like Active-Bidder will be essential in ensuring that these systems serve to reduce, rather than increase, the cognitive load and ``informational surprise'' of the human users they interact with.